\section{第2章 导数与微分}
\subsection{导数的基本概念}
\begin{exercise}{151}{基础·导数定义}
用导数定义求 $f(x)=1+2x+3x^2$ 在 $x=1$ 处的导数。
\end{exercise}
\begin{exercise}{152}{基础·零点处求导}
不展开多项式，求 $f(x)=x(2x+1)(3x+2)(4x+3)$ 的 $f'(0)$。
\end{exercise}
\begin{exercise}{153}{常规·定义求导与切线}
用导数定义求 $f(x)=x(x-2)\sqrt{\frac{3x^2+1}{2x^2+1}}$ 的 $f'(0)$，并求曲线在原点的切线。
\end{exercise}
\begin{exercise}{154}{基础·补定义函数的导数}
设 $f(x)=\frac{\sin^2x}{2x}$（$x\ne0$），$f(0)=0$。用定义求 $f'(0)$。
\end{exercise}
\begin{exercise}{155}{常规·连续因子与求导}
设 $f(x)=(x^2-1)\varphi(x-1)$，$\varphi$ 在 $0$ 处连续且 $\varphi(0)=3$。用定义求 $f'(1)$。
\end{exercise}
\begin{exercise}{156}{常规·平移差商}
设 $f$ 在 $0$ 处可导，$f'(0)=\frac{1}{2}$，且 $f(2+x)=2f(x)$ 对所有实数 $x$ 成立。用定义求 $f'(2)$。
\end{exercise}
\begin{exercise}{157}{基础·差商极限}
设 $f$ 在 $0$ 处连续，且 $\lim_{x\to0}\frac{f(x)+3}{x}=2$。求 $f(0)$ 和 $f'(0)$。
\end{exercise}
\begin{exercise}{158}{基础·增量缩放}
设 $f'(x_0)=2$。求 $\lim_{h\to0}\frac{f(x_0+5h)-f(x_0)}{h}$ 及 $\lim_{h\to0}\frac{f(x_0)-f(x_0-3h)}{h}$。
\end{exercise}
\begin{exercise}{159}{常规·双移位差商}
设 $f'(x_0)=3$。求 $\lim_{h\to0}\frac{f(x_0+4h)-f(x_0+2h)}{h}$。提示：在分子中加减 $f(x_0)$。
\end{exercise}
\begin{exercise}{160}{常规·二次增量差商}
设 $f$ 在 $3$ 处可导，且 $\lim_{x\to0}\frac{f(3)-f(3-3x^2)}{\ln(1+x^2)}=3$。求 $f'(3)$。
\end{exercise}
\begin{exercise}{161}{基础·已知导数求极限}
设 $f(0)=0$、$f'(0)=1$。求 $\lim_{x\to0}\frac{f(3x)-f(-5x)}{x}$。
\end{exercise}
\begin{exercise}{162}{基础·导数定义换元}
设 $f(0)=0$、$f'(0)=2$。求 $\lim_{h\to0}\frac{f(h^3)}{h^3}$ 和 $\lim_{h\to0}\frac{f(h^2)}{h^2}$。
\end{exercise}
\begin{exercise}{163}{常规·复合增量极限}
设 $f$ 在 $1$ 处可导，$f(1)=0$、$f'(1)=2$。求 $\lim_{x\to0}\frac{f(\cos x)}{x^2}$。
\end{exercise}
\begin{exercise}{164}{基础·左右导数}
分别用左右差商判断 $|x|$、$x^2$ 在 $0$ 处是否可导，可导时求导数。
\end{exercise}
\begin{exercise}{165}{基础·分段函数可导性}
设 $f(x)=\cos x$（$x>0$），$f(x)=x^2+1$（$x\le0$）。求 $f'_-(0)$、$f'_+(0)$，判断可导性。
\end{exercise}
\begin{exercise}{166}{常规·指数分段函数}
设 $f(x)=x(1+e^{\frac{1}{x}})$（$x\ne0$），$f(0)=0$。判断 $f$ 在 $0$ 处的连续性，并求左导数（若存在）。
\end{exercise}
\begin{exercise}{167}{基础·绝对值尖点}
求 $f(x)=x^2|x-1|$ 在 $1$ 处的左右导数，判断是否可导。
\end{exercise}
\begin{exercise}{168}{常规·零点消除尖点}
用定义求 $f(x)=(x^3-1)|x-1|$ 在 $1$ 处的导数。
\end{exercise}
\begin{exercise}{169}{常规·连续因子差商}
设 $\varphi$ 在 $1$ 处连续且 $\varphi(1)=0$，$f(x)=|x-1|\varphi(x)$。用 $|\varphi(x)|\to0$ 证明 $f'(1)=0$。
\end{exercise}
\begin{exercise}{170}{常规·绝对值复合}
求 $f(x)=|x^2-1|$ 的不可导点，并计算这些点的左右导数。
\end{exercise}
\begin{exercise}{171}{基础·切线与法线}
求曲线 $y=2x+\cos x$ 在 $(0,1)$ 处的切线与法线方程。
\end{exercise}
\begin{exercise}{172}{常规·切线垂直}
求实数 $x_0$，使曲线 $y=x^2$ 与 $y=x^3$ 在横坐标同为 $x_0$ 的点处的切线互相垂直。
\end{exercise}
\begin{exercise}{173}{常规·指定方向的切线}
求与直线 $2x-6y+1=0$ 垂直且与曲线 $y=x^3+3x^2-5$ 相切的直线方程。
\end{exercise}
\begin{exercise}{174}{常规·切线与面积}
求双曲线 $xy=1$ 在 $(1,1)$ 处的切线，并求该切线与坐标轴围成的三角形面积。
\end{exercise}
\begin{exercise}{175}{常规·公共切线}
求 $a,b$，使 $y=1+\ln x$ 与 $y=x^2+ax+b$ 在 $x=1$ 处具有同一条切线。
\end{exercise}
\begin{exercise}{176}{常规·偶函数导数}
设 $f$ 为可导偶函数且 $f'(3)=-\frac{1}{2}$。求曲线在 $(-3,f(-3))$ 处的切线斜率。可使用偶函数的导函数为奇函数。
\end{exercise}
\begin{exercise}{177}{基础·可导与连续}
设 $f$ 在 $2$ 处可导且 $\lim_{x\to2}f(x)=5$。求 $f(2)$，并说明依据。
\end{exercise}
\begin{exercise}{178}{常规·振荡函数可导性}
设 $f(x)=x^2\cos(\frac{1}{x})$（$x\ne0$），$f(0)=0$。用定义判断其在 $0$ 处的连续性与可导性。
\end{exercise}
\begin{exercise}{179}{常规·分段拼接}
设 $f(x)=xe^x$（$x\ge0$），$f(x)=x^3+x$（$x<0$）。判断其在 $0$ 处是否可导，并求 $f'(0)$。
\end{exercise}
\begin{exercise}{180}{基础·先连续再可导}
设 $f(x)=x^3+\ln x$（$x\ge1$），$f(x)=x^3\ln x$（$0<x<1$）。判断 $f$ 在 $1$ 处是否连续、可导。
\end{exercise}
\begin{exercise}{181}{常规·可导拼接参数}
求 $a,b$，使 $f(x)=x\sin\pi x$（$x>1$），$f(x)=ax^2+b$（$x\le1$）在 $1$ 处可导。
\end{exercise}
\begin{exercise}{182}{常规·有界因子与可导}
设 $\varphi$ 在 $0$ 附近有界，$f(x)=x^2\varphi(x)$ 且 $f(0)=0$。用定义求 $f'(0)$。
\end{exercise}
\begin{exercise}{183}{基础·导数定义变式}
用定义求 $f(x)=\sin2x$ 的 $f'(0)$。
\end{exercise}
\begin{exercise}{184}{常规·分段可导综合}
设 $f(x)=ax+b$（$x<1$），$f(x)=x^2$（$x\ge1$）。求使 $f$ 在 $1$ 处可导的 $a,b$，并写出切线方程。
\end{exercise}
\subsection{函数的求导法则}
\begin{exercise}{185}{基础·基本求导}
求 $f(x)=3x^2+2e^x-5\ln x+\sin1$ 的导数及 $f'(1)$。
\end{exercise}
\begin{exercise}{186}{基础·乘积法则}
求 $f(x)=x\cos x+3x^2$ 的导数及 $f'(\frac{\pi}{2})$。
\end{exercise}
\begin{exercise}{187}{基础·指数三角乘积}
求 $y=e^x(\cos x-2\sin x)$ 的导数。
\end{exercise}
\begin{exercise}{188}{基础·乘积与复合求导}
求 $y=(1+x^2)\arctan x-x\ln x$（$x>0$）的导数。
\end{exercise}
\begin{exercise}{189}{基础·商式求导}
求 $y=\frac{\sqrt x+xe^x-3x^3\sin x}{x^2}$（$x>0$）的导数。提示：可先拆分为三项。
\end{exercise}
\begin{exercise}{190}{基础·商法则}
求 $y=\frac{\sin x}{x-\cos x}$ 的导数，注明分母限制。
\end{exercise}
\begin{exercise}{191}{基础·对数有理式}
求 $y=\frac{2\ln x+x^3}{3\ln x+x^2}$ 的导数及 $y'(1)$。
\end{exercise}
\begin{exercise}{192}{基础·乘商求导}
求 $y=\frac{e^x(x^2+1)}{x+1}$（$x\ne-1$）的导数。
\end{exercise}
\begin{exercise}{193}{常规·换元与求导}
已知 $f(x-1)=x^2+1$ 对所有实数 $x$ 成立。求 $f(t)$ 和 $f'(t)$。
\end{exercise}
\begin{exercise}{194}{常规·分段导函数}
设 $f(x)=\frac{x}{2x-1}$（$x<0$），$f(x)=e^x\sin x$（$x\ge0$）。分别求 $x<0$、$x>0$ 时的导数，并判断 $0$ 处是否可导。
\end{exercise}
\begin{exercise}{195}{基础·反函数导数}
设 $f$ 的可导反函数为 $\varphi$，已知 $f(1)=2$、$f(2)=3$、$f'(1)=4$、$f'(2)=5$。求 $\varphi'(2)$、$\varphi'(3)$。
\end{exercise}
\begin{exercise}{196}{常规·反函数定点导数}
设 $f(x)=3x^3+2x+1$ 的反函数为 $\varphi$。求 $\varphi'(1)$。提示：先求 $f(x)=1$ 的解。
\end{exercise}
\begin{exercise}{197}{基础·链式法则}
求 $y=(1-x+2x^3)^4$ 的导数及 $y'(0)$。
\end{exercise}
\begin{exercise}{198}{基础·多层复合求导}
求 $y=\ln(\cos(e^{-x}))$ 在 $x=0$ 附近的导数及 $y'(0)$。
\end{exercise}
\begin{exercise}{199}{常规·反三角复合求导}
求 $f(x)=\arcsin\sqrt{\frac{1-x}{1+x}}$ 在 $0<x<1$ 内的导数及 $f'(\frac{1}{2})$。
\end{exercise}
\begin{exercise}{200}{常规·根式商求导}
求 $y=\frac{\sqrt{x+1}-\sqrt{x-1}}{\sqrt{x+1}+\sqrt{x-1}}$（$x>1$）的导数。提示：可先化简为 $x-\sqrt{x^2-1}$。
\end{exercise}
\begin{exercise}{201}{基础·指数与幂指求导}
求 $y=3^x+x^3+x^{2x}$（$x>0$）的导数及 $y'(1)$。
\end{exercise}
\begin{exercise}{202}{基础·三角复合求导}
求 $y=\cos3x\,\sin^3x$ 的导数及 $y'(\frac{\pi}{6})$。
\end{exercise}
\begin{exercise}{203}{基础·反三角与根式求导}
求 $y=x\sqrt{1-x^2}-\arccos x$ 在 $-1<x<1$ 内的导数。
\end{exercise}
\begin{exercise}{204}{常规·根式对数求导}
求 $y=x^2+x\sqrt{x^2+1}+\ln(x+\sqrt{x^2+1})$ 的导数并化简。
\end{exercise}
\begin{exercise}{205}{基础·对数三角复合}
求 $y=x[\cos(\ln x)+\sin(\ln x)]$（$x>0$）的导数。
\end{exercise}
\begin{exercise}{206}{基础·反余切复合}
求 $y=e^{-x}\operatorname{arccot}(e^x)$ 的导数，其中 $\operatorname{arccot}\in(0,\pi)$。
\end{exercise}
\begin{exercise}{207}{常规·先求极限再求导}
对固定实数 $x$，设 $f(x)=\lim_{n\to\infty}x(1-\frac{3x}{n})^n$，只考虑底数为正的充分大的 $n$。求 $f(x)$ 和 $f'(x)$。
\end{exercise}
\begin{exercise}{208}{基础·抽象函数链式法则}
设 $f$ 可导，求 $y=f(3^x+x^2)$ 的导数；若 $f'(1)=2$，求 $y'(0)$。
\end{exercise}
\begin{exercise}{209}{常规·复合函数取点}
已知 $f$ 可导且 $f'(1)=-2$，$g(x)=f(2x^3-3x)$。求 $g'(-1)$。
\end{exercise}
\begin{exercise}{210}{基础·复合求导变式}
求 $y=\ln(1+x^2)$ 的导数及在 $x=1$ 处的切线。
\end{exercise}
\begin{exercise}{211}{常规·反函数复合求导}
设 $f(x)=x+e^x$ 的反函数为 $\varphi$，$G(x)=\varphi(e^x)$。求 $\varphi'(1)$ 和 $G'(0)$。
\end{exercise}
\subsection{高阶导数}
\begin{exercise}{212}{基础·二阶导数}
分别求 $y=\sin2x$ 和 $y=x\arctan x$ 的二阶导数。
\end{exercise}
\begin{exercise}{213}{基础·指数三角二阶导数}
设 $f(x)=e^{3x}\sin2x$。求 $f''(x)$ 和 $f''(\frac{\pi}{2})$。
\end{exercise}
\begin{exercise}{214}{基础·对数函数二阶导数}
设 $y=\frac12\ln(1+e^{2x})-x$。求 $y'(0)$ 和 $y''(0)$。
\end{exercise}
\begin{exercise}{215}{基础·反正弦二阶导数}
设 $y=\arcsin(\frac{x}{3})$。求 $y''(x)$（$|x|<3$）及 $y''(\sqrt5)$。
\end{exercise}
\begin{exercise}{216}{常规·高阶导数通式}
设 $y=x\ln x$（$x>0$）。求 $y''$、$y'''$，并归纳写出 $y^{(n)}$（$n\ge2$）的通式。
\end{exercise}
\begin{exercise}{217}{基础·三角高阶导数}
设 $y=(\sin(\frac{x}{2})+\cos(\frac{x}{2}))^2$。先化简，再求 $y^{(10)}$。
\end{exercise}
\begin{exercise}{218}{常规·莱布尼茨公式}
设 $f(x)=xe^{-x}$。用莱布尼茨公式求 $f^{(n)}(x)$（$n\ge1$）。
\end{exercise}
\begin{exercise}{219}{基础·多项式高阶导数}
设 $y=(x^2+1)^5(x^5+1)^2$。不完全展开，求 $y^{(20)}$ 和 $y^{(21)}$。
\end{exercise}
\begin{exercise}{220}{常规·降幂求导}
设 $f(x)=\sin^4x+\cos^4x$。利用降幂公式求 $f^{(4)}(x)$。
\end{exercise}
\begin{exercise}{221}{常规·对数高阶导数}
设 $f(x)=\ln(2-3x)$（$x<\frac{2}{3}$）。求 $f'$、$f''$、$f'''$，并写出 $f^{(n)}$ 的通式。
\end{exercise}
\begin{exercise}{222}{常规·导数关系再求导}
设 $f'=f^2$ 且 $f(1)=2$。求 $f''(1)$。已知 $f$ 在 $1$ 附近可导。
\end{exercise}
\begin{exercise}{223}{基础·高阶导数变式}
求 $y=e^{2x}$ 的四阶导数及 $y^{(4)}(0)$。
\end{exercise}
\begin{exercise}{224}{常规·乘积高阶导数综合}
用莱布尼茨公式求 $f(x)=x^2e^x$ 的 $f^{(4)}(x)$，并求 $f^{(4)}(0)$。
\end{exercise}
\subsection{隐函数及由参数方程所确定的函数的导数}
\begin{exercise}{225}{基础·隐函数求导}
由方程 $x^2+y^2=6xy$ 求 $\frac{\dd y}{\dd x}$，注明公式分母非零的条件。
\end{exercise}
\begin{exercise}{226}{常规·隐函数定点导数}
方程 $xy-\sin(\pi y^2)=0$ 确定过 $(0,1)$ 的可导分支。求 $y'(0)$。
\end{exercise}
\begin{exercise}{227}{常规·隐曲线切线}
求曲线 $\tan(x+y+\frac{\pi}{4})=y$ 在 $(-1,1)$ 处的切线与法线。
\end{exercise}
\begin{exercise}{228}{常规·隐函数分支}
求曲线 $y^2+2e^{xy}=3y$ 在 $(0,1)$ 处的切线方程。
\end{exercise}
\begin{exercise}{229}{常规·抽象隐函数求导}
设 $xf(y)=e^y\ln y$ 确定过 $(0,1)$ 的可导分支，$f$ 可导且 $f(1)=1$。求 $y'(0)$。
\end{exercise}
\begin{exercise}{230}{常规·公共切线参数}
曲线 $y=x^2+ax+b$ 与 $2y=xy^2-3$ 在 $(1,-1)$ 处有公共切线。求 $a,b$。
\end{exercise}
\begin{exercise}{231}{常规·隐函数二阶导数}
方程 $x^2+y^2=2$ 确定过 $(1,1)$ 的分支。求 $y'(1)$ 和 $y''(1)$。
\end{exercise}
\begin{exercise}{232}{常规·指数隐函数二阶导数}
设 $e^y+xy=e+1$ 确定过 $(1,1)$ 的局部分支。求 $y'(1)$、$y''(1)$。
\end{exercise}
\begin{exercise}{233}{常规·含导数的函数关系}
设 $y(x)$ 二阶可导，经过 $(1,-1)$，满足 $x^2+y^2+(y')^3=1$。求 $y'(1)$、$y''(1)$。
\end{exercise}
\begin{exercise}{234}{基础·对数求导}
用对数求导法求 $y=x^{\ln x}$（$x>0$）的导数。
\end{exercise}
\begin{exercise}{235}{基础·变底变指数求导}
设 $y=(1+\sin x)^{x^2}$，其中 $1+\sin x>0$。求 $y'$ 和 $y'(0)$。
\end{exercise}
\begin{exercise}{236}{常规·幂指函数乘积}
求 $y=x^{\sin x}\sqrt{e^x-1}$（$x>0$）的导数。
\end{exercise}
\begin{exercise}{237}{常规·复合幂指函数}
求 $y=(1+2x)^{\frac{1}{x}}e^{-x}$（$x>-\frac{1}{2}$、$x\ne0$）的导数。
\end{exercise}
\begin{exercise}{238}{常规·多因子对数求导}
用对数求导法求 $y=\frac{(x+1)^2\sqrt{2x+2}}{(x-1)^3}$（$x>1$）的导数。
\end{exercise}
\begin{exercise}{239}{基础·幂指函数}
求 $y=x^x$（$x>0$）的导数及驻点。
\end{exercise}
\begin{exercise}{240}{常规·幂指二阶导数}
求 $y=x^{2x}$（$x>0$）的二阶导数。
\end{exercise}
\begin{exercise}{241}{常规·对数化隐函数}
在 $x,y>0$ 且 $1+\ln y\ne0$ 的可导分支上，方程 $x^{\frac{1}{y}}=y^{\frac{1}{x}}$ 确定 $y=y(x)$。取对数化简后求 $y'$。
\end{exercise}
\begin{exercise}{242}{基础·参数方程求导}
求参数曲线 $x=\frac{2t}{1+t^2}$、$y=\frac{2t^2}{1+t^2}$ 的 $\frac{\dd y}{\dd x}$，注明 $\frac{\dd x}{\dd t}\ne0$ 的限制。
\end{exercise}
\begin{exercise}{243}{基础·参数定点导数}
设 $x=t^2e^{2t}$、$y=(t^2+1)e^t$。求 $t=2$ 时的 $\frac{\dd y}{\dd x}$。
\end{exercise}
\begin{exercise}{244}{常规·参数曲线切线}
求参数曲线 $x=\sqrt{t^2+3}$、$y=\ln(t+\sqrt{t^2+3})$ 在 $t=1$ 对应点处的切线。
\end{exercise}
\begin{exercise}{245}{常规·参数曲线法线}
求参数曲线 $x=e^t\sin2t$、$y=e^t\cos t$ 在 $t=0$ 对应点处的切线和法线。
\end{exercise}
\begin{exercise}{246}{常规·参数二阶导数}
设 $x=\cos t$、$y=t\cos t-\sin t$，$0<t<\pi$。求 $\frac{\dd y}{\dd x}$ 与 $\frac{\dd^2y}{\dd x^2}$。
\end{exercise}
\begin{exercise}{247}{常规·参数二阶导数取值}
设 $x=\arctan t$、$y=3t+t^3$。求 $t=1$ 时的 $\frac{\dd^2y}{\dd x^2}$。
\end{exercise}
\begin{exercise}{248}{常规·参数导数综合}
设 $x=e^t+t$、$y=\sin t$。求 $t=0$ 时的 $\frac{\dd y}{\dd x}$ 与 $\frac{\dd^2y}{\dd x^2}$。
\end{exercise}
\begin{exercise}{249}{基础·隐函数求导变式}
求曲线 $x^2+xy+y^2=3$ 在 $(1,1)$ 处的导数与切线。
\end{exercise}
\begin{exercise}{250}{常规·参数曲线小综合}
设 $x=t^2+1$、$y=t^3$。求 $t=1$ 对应点的切线及该点处的二阶导数。
\end{exercise}
\subsection{函数的微分}
\begin{exercise}{251}{基础·微分与增量}
设 $y=x^3+x$。在 $x=3$、$\Delta x=0.1$ 时计算 $\dd y$ 和 $\Delta y$。
\end{exercise}
\begin{exercise}{252}{基础·由增量求微分}
已知 $f$ 的增量为 $\Delta y=\frac{\Delta x}{\sqrt{1+x^3}}+o(\Delta x)$（$x>-1$）。求 $f'(2)$ 及 $x=2$ 处的微分。
\end{exercise}
\begin{exercise}{253}{常规·微分的线性主部}
设 $f$ 在 $x_0$ 处可导且 $f'(x_0)=1$。写出 $\dd y$，并求 $\lim_{\Delta x\to0}\frac{\Delta y}{\Delta x}$。
\end{exercise}
\begin{exercise}{254}{基础·复合函数微分}
求 $y=e^{\sin^2x}$ 和 $y=\arcsin(\frac{x}{2})$（$|x|<2$）的微分。
\end{exercise}
\begin{exercise}{255}{基础·定点微分}
分别求 $y=\sqrt2+\sin x$ 在 $x=0$ 处、$y=\ln(1+2^{-2x})$ 在 $x=1$ 处的微分。
\end{exercise}
\begin{exercise}{256}{常规·隐函数微分}
设 $2y+\sin(x+y)=0$ 确定可导函数 $y=y(x)$。求 $\dd y$ 及原点处的微分。
\end{exercise}
\begin{exercise}{257}{常规·隐函数定点微分}
方程 $e^y+xy-e=0$ 在 $(0,1)$ 附近确定分支。求 $x=0$ 处的微分，并估计 $x$ 增加 $0.01$ 时 $y$ 的增量。
\end{exercise}
\begin{exercise}{258}{常规·补定义点微分}
设 $f(x)=\frac{\cos x-1}{x}$（$x\ne0$），$f(0)=0$。用导数定义求 $x=0$ 处的微分。
\end{exercise}
\begin{exercise}{259}{基础·幂指函数微分}
求 $y=(1+e^x)^x$ 的微分及 $x=0$ 处的微分。
\end{exercise}
\begin{exercise}{260}{基础·由导数求函数}
求导数分别为 $x^2$、$\cos3x$、$e^x$ 的全部函数，各式均在实轴上考虑。
\end{exercise}
\begin{exercise}{261}{常规·夹逼与可微性}
设 $f$ 在 $0$ 附近有定义且 $|f(x)|\le x^2$。用导数定义证明 $f'(0)=0$，并写出该点的微分。
\end{exercise}
\begin{exercise}{262}{常规·定点导数计算}
分别求 $p(x)=x(x-1)(x-2)(x-3)$ 的 $p'(3)$ 和 $y=\frac{(x^2+1)\sqrt{3x+1}}{(x+1)^2e^{2x}}$ 的 $y'(1)$。
\end{exercise}
\begin{exercise}{263}{常规·差商反求导数}
设 $f$ 在 $x_0$ 处可导，$f(x_0)=0$，且 $\lim_{h\to0}\frac{f(x_0+3h)-3f(x_0-h)}{h}=1$。求 $f'(x_0)$。
\end{exercise}
\begin{exercise}{264}{常规·可导信息与幂指极限}
设 $f$ 在 $0$ 处可导，$f(0)=1$、$f'(0)=4$。求 $\lim_{x\to0}[f(\sin x)]^{\frac{1}{2x}}$。提示：取对数并使用导数定义。
\end{exercise}
\begin{exercise}{265}{常规·切线截距}
求曲线 $y=\frac{1}{1+x^2}$ 在 $(1,\frac{1}{2})$ 处的切线及它与 $x$ 轴的交点。
\end{exercise}
\begin{exercise}{266}{常规·复合差商与法线}
设 $f$ 连续可导，且 $\lim_{x\to1}\frac{f(3x-1)+3}{x-1}=3$。求 $f(2)$、$f'(2)$ 及该点的法线方程。
\end{exercise}
\begin{exercise}{267}{常规·分段函数微分}
设 $f(x)=x^2$（$|x|\le1$），$f(x)=1$（$|x|>1$）。判断 $x=1$ 处是否可微，并求 $x=0$ 处的微分。
\end{exercise}
\begin{exercise}{268}{常规·连续补值}
设 $f$ 在 $1$ 处可导，$f(1)=0$、$f'(1)=1$。令 $g(x)=\frac{f(e^{2x})+\ln(1-4x)}{x}$（$x\ne0$），$g(0)=b$，在 $0$ 附近考虑。求使 $g$ 连续的 $b$。
\end{exercise}
\begin{exercise}{269}{常规·分段可导参数}
设 $f(x)=x^2$（$x\ge1$），$f(x)=3ax+b$（$x<1$）。求使 $f$ 在 $1$ 处可导的 $a,b$。
\end{exercise}
\begin{exercise}{270}{基础·互补复合求导}
设 $f$ 可导，$y=f(\sin^2x)+f(\cos^2x)$。求 $y'$ 与 $y'(0)$。
\end{exercise}
\begin{exercise}{271}{常规·两层抽象复合}
设 $f$ 可导，$y=f(f(1+3x^2))$。求 $y'$ 与 $y'(0)$。
\end{exercise}
\begin{exercise}{272}{常规·公共切线与复合导数}
曲线 $y=f(x)$ 与 $y=\frac{3+\ln x}{x}$ 在 $x=1$ 处切线相同。设 $F(x)=[f(e^x)]^2$，求 $F'(0)$。
\end{exercise}
\begin{exercise}{273}{常规·二阶导数与极限}
设 $f$ 在 $0$ 附近二阶连续可导，$f(0)=f'(0)=0$、$f''(0)=2$。求 $\lim_{x\to0}\frac{e^{f(x)}-1}{1-\cos x}$。提示：使用二阶泰勒展开。
\end{exercise}
\begin{exercise}{274}{常规·复合二阶导数}
设 $f$ 二阶可导且 $f'(0)=3$，$y=f(x^2)$。求 $y''(0)$。
\end{exercise}
\begin{exercise}{275}{常规·导数关系计算}
设 $f'=e^f$ 且 $f(0)=0$。求 $f'(0)$ 和 $f''(0)$。
\end{exercise}
\begin{exercise}{276}{常规·莱布尼茨定点求导}
设 $f(x)=(x+1)^3e^{-x^2}$。用莱布尼茨公式求 $f'''(-1)$。
\end{exercise}
\begin{exercise}{277}{常规·反函数二阶导数}
设 $x=\varphi(y)$ 是 $y=x+\ln x$（$x>0$）的反函数。求 $y=1$ 时的 $\frac{\dd x}{\dd y}$ 与 $\frac{\dd^2x}{\dd y^2}$。
\end{exercise}
\begin{exercise}{278}{常规·振荡函数定点求导}
设 $f(x)=x^3\sin(\frac{1}{x})+\sin x+1$（$x\ne0$），$f(0)=1$。用定义求 $f'(0)$，并求 $x\ne0$ 时的导数。
\end{exercise}
\begin{exercise}{279}{常规·隐函数与局部极限}
方程 $\ln(x^2+y)=x^2y+\sin x$ 在 $(0,1)$ 附近确定分支 $y=f(x)$。求 $f'(0)$ 及 $\lim_{h\to0}\frac{f(h)-1}{h}$。
\end{exercise}
\begin{exercise}{280}{常规·隐函数复合求导}
设 $f$ 可导且 $f'(1)=1$，$xy-e^y+e^x=0$ 确定过原点的可导分支。令 $z=f(e^{y(x)})$，求 $z'(0)$。
\end{exercise}
\begin{exercise}{281}{基础·参数定点微分}
设 $x=2t^2+3t+1$、$y=e^t\sin t+2$。求 $t=0$ 时的 $\frac{\dd y}{\dd x}$ 及该点的微分。
\end{exercise}
\begin{exercise}{282}{常规·复合参数求导}
设 $x=t-\frac12\ln(1+e^{2t})$、$y=\arctan(e^t)$。求 $t=0$ 时的 $\frac{\dd y}{\dd x}$。
\end{exercise}
\begin{exercise}{283}{常规·极坐标曲线切线}
求螺线 $\rho=\theta$ 在 $\theta=\frac{\pi}{2}$ 处的切线方程。提示：使用 $x=\theta\cos\theta$、$y=\theta\sin\theta$。
\end{exercise}
\begin{exercise}{284}{基础·由微分求导数}
设 $y=f(e^{2x})$，在 $x=0$ 处有 $\dd y=\dd x$。求 $f'(1)$。
\end{exercise}
\begin{exercise}{285}{常规·复合定点微分}
设 $f'(x)=\arctan(\frac{1}{x})$（$x\ne0$），$y=f(\frac{2x-1}{x+1})$。求 $x=2$ 处的微分。
\end{exercise}
\begin{exercise}{286}{常规·参数函数微分}
设 $f$ 可导且 $f'(0)\ne0$，$x=f(\ln t)$、$y=f(t^2-1)$（$t>0$）。求 $t=1$ 时的 $\dd y$，分别用 $\dd x$ 与 $\dd t$ 表示。
\end{exercise}
\begin{exercise}{287}{基础·微分近似}
用微分估算 $\sqrt{4.04}$。
\end{exercise}
\begin{exercise}{288}{常规·隐函数微分综合}
设 $ye^y=x$ 在 $(e,1)$ 附近确定 $y(x)$。求该点的微分，并用它估算 $x=e+0.02$ 时的 $y$。
\end{exercise}
